\documentclass[11pt]{amsart}
\usepackage[a4paper,margin=25mm]{geometry}
\usepackage[T1]{fontenc}
\usepackage{lmodern,amsmath,amssymb,amsthm,mathtools,microtype}
\usepackage[hidelinks]{hyperref}
\hypersetup{pdftitle={Uniform condensation windows for zero-radius random allocations},pdfauthor={Yitzchak Shmalo},pdfsubject={Research draft with five documented self-audit passes},pdfkeywords={condensation, random allocations, simply generated trees}}
\usepackage{enumitem,booktabs}
\setlist{nosep,leftmargin=*}
\newtheorem{theorem}{Theorem}[section]
\newtheorem{lemma}[theorem]{Lemma}
\newtheorem{corollary}[theorem]{Corollary}
\newtheorem{proposition}[theorem]{Proposition}
\theoremstyle{remark}\newtheorem{remark}[theorem]{Remark}
\newcommand{\Pp}{\mathbb P}
\newcommand{\E}{\mathbb E}
\newcommand{\N}{\mathbb N}
\newcommand{\supp}{\operatorname{supp}}
\newcommand{\one}{\mathbf 1}
\newcommand{\fall}[2]{(#1)_{\underline{#2}}}
\title[Uniform condensation windows]{Uniform condensation windows for zero-radius random allocations}
\author{Yitzchak Shmalo}
\address{Einstein Institute of Mathematics, The Hebrew University of Jerusalem, Jerusalem, Israel}
\date{September 17, 2026}
\subjclass[2020]{60C05, 60J80, 05C05}
\keywords{Random allocation, condensation, simply generated tree, zero-radius weights, largest degree}
\begin{document}
\begin{abstract}
For any fixed nonnegative weight sequence with positive constant term and
zero radius of convergence, we prove the existence of arbitrarily large
intervals of total masses on which a single box carries almost all the mass,
uniformly over the number of boxes at densities in a compact subset of
$(0,\infty)$. Consequently every admissible full mass path of positive finite
limiting density has a subsequence exhibiting complete condensation. This
answers the zero-radius question in Janson's survey (Problem 19.41 in the
published version, Problem 18.40 in the preprint).
The argument preserves the small-box partition factor by an exact
occupation-count insertion. A deterministic record-margin estimate then
provides a uniform finite-size probability bound. We also show that
multiplicative density of a sampled index set is necessary and sufficient
for a universal subsequence-condensation statement.
\end{abstract}

\maketitle

\section{Introduction}\label{sec:intro}
Weighted random allocations provide a common framework for occupancy models
and degree sequences of simply generated trees. The probability of an
allocation is proportional to the product of the weights of its box sizes.
When the total number of balls is proportional to the number of boxes, a
basic question is whether one box carries a macroscopic fraction of the
mass. For irregular weights this question is sensitive to the sizes at
which the model is observed. Janson's survey \cite{Janson} develops the
allocation--tree correspondence and formulates several questions about
such condensation.

We consider fixed weights $w_j\ge0$ with $w_0>0$ for which the formal
generating function $\Phi(z)=\sum_{j\ge0}w_jz^j$ has radius zero. This is
precisely the nontrivial $\nu=0$ regime of the survey. No assumption is made
on the ratios $w_{j+1}/w_j$ or on the support. In particular, a large weight
may be followed by an arbitrarily long interval of zero weights. A single
probability law obtained by a fixed positive exponential tilt is unavailable:
$\Phi(z)=\infty$ for every $z>0$. We instead use finite truncations and
record-dependent tilts, always retaining the original allocation
partition function in the denominator.

The question addressed here is whether the largest box can be $o(n)$ in
probability throughout a prescribed sequence $m_n/n\to\lambda>0$.
Janson's Problem~19.41 \cite[p.~228]{Janson}, numbered Problem~18.40
in the preprint \cite{JansonPreprint}, asks this in the zero-radius regime.
The answer given below is negative. More strongly, for every such path
there is a subsequence along which a single box carries a fraction tending
to one of the total mass. The main theorem is uniform over entire intervals
of masses and over box counts in any fixed compact positive-density range.
The intervals are chosen from the weights, rather than separately for each
path. Arbitrary jumps and oscillations in $m_n$ are permitted.

There are two obstructions to a direct comparison with an all-in-one-box
configuration. First, the target mass need not belong to the support of the
weights. Second, the total contribution of configurations consisting of a
giant box and many small boxes need not be comparable to the weight of a
single configuration. We resolve the first obstruction by a supported
representation of the residual mass, and the second by an exact insertion
inequality that keeps the whole small-box partition factor. Its factorial
count compensates for the number of inserted large residual pieces.
A geometric record decomposition bounds their weight loss. Summing that
loss requires distinguishing scales at which the record increases from
possibly long intervals on which it is constant.

The conclusion is genuinely subsequential. Sparse factorial support gives
examples with a sublinear largest box on certain subsequences
\cite[Example~19.39]{Janson}. In a different direction, Janson, Jonsson and
Stef\'ansson \cite{JJS} prove an infinite-star local limit under a
successive-weight-ratio hypothesis, and stronger size estimates for
explicit factorial families. Their local convergence and special-family
results differ from the uniform subsequence statement here, which imposes
only zero radius. We do not address the positive-boundary-mean question
(Problem~19.40 of \cite{Janson}, Problem~18.39 of \cite{JansonPreprint}).

We also identify the exact sampling condition for the universal statement.
An increasing sampled index set admits a near-total-condensation
subsequence for every zero-radius law and every admissible positive-density
path if and only if the ratio of successive sampled indices tends to one.
The converse is a deterministic support construction. For simply generated
plane trees, the allocation theorem and the known size-biased root identity
imply that, along a subsequence, the root has almost all vertices as children.

\paragraph{Organization.}
Section~\ref{sec:results} states the main results. Sections~\ref{sec:insertion}
and \ref{sec:records} give the insertion inequality and the geometric
summation bounds. Section~\ref{sec:residual} constructs supported residuals,
and Section~\ref{sec:margin} gives a deterministic favorable-margin
criterion. Section~\ref{sec:windows} proves the uniform theorem.
Sections~\ref{sec:paths} and \ref{sec:sharp} treat paths, trees, and sampling
sharpness. Section~\ref{sec:finite} records exact finite computations; none
of these computations is used as a substitute for the proofs.


\section{Model and results}\label{sec:results}
Let $w=(w_j)_{j\ge0}$ be fixed, with $w_0>0$, $0\le w_j<\infty$, and
\begin{equation}\label{eq:radius}
 \rho=0,\qquad\text{equivalently}\qquad
 \limsup_{j\to\infty}w_j^{1/j}=\infty.
\end{equation}
For nonnegative integers $m,n$ put
\[
 Z(m,n)=\sum_{y_1+\cdots+y_n=m}\prod_{v=1}^n w_{y_v}.
\]
The pair $(m,n)$ is \emph{admissible} if $Z(m,n)>0$. On such a pair,
\[
 \Pp_w(Y=y)=Z(m,n)^{-1}\prod_{v=1}^n w_{y_v},
 \qquad Y_{(1)}\ge Y_{(2)}\ge\cdots
\]
denote the ordered occupancies. The weights are independent of $m,n$.

\begin{theorem}[Uniform condensation windows]\label{thm:windows}
Assume \eqref{eq:radius}. Fix $\varepsilon\in(0,1)$ and
$0<\alpha<\beta<\infty$. There are $\delta>0$ and supported integers
$k_r\uparrow\infty$ such that
\begin{equation}\label{eq:uniform}
 \sup_{\substack{(1+3\delta)k_r\le m\le(1+5\delta)k_r\\
                  \alpha\le m/n\le\beta,\ Z(m,n)>0}}
 \Pp_w\bigl(Y_{(1)}(m,n)\le(1-\varepsilon)m\bigr)
 \longrightarrow0.
\end{equation}
The supremum is over integer pairs. The windows contain admissible pairs
for all sufficiently large $r$; thus the conclusion is not vacuous.
\end{theorem}

\begin{corollary}[Arbitrary full paths and dense sampled paths]\label{cor:paths}
Let $\mathcal N=\{N_1<N_2<\cdots\}$ satisfy $N_{j+1}/N_j\to1$; in particular
one may take $\mathcal N=\N$. Suppose $(m_n,n)$ is eventually admissible on
$\mathcal N$ and $m_n/n\to\lambda\in(0,\infty)$. Then there are
$n_r\in\mathcal N$, increasing to infinity, such that
\begin{equation}\label{eq:conclusion}
 \frac{Y_{(1)}(m_{n_r},n_r)}{m_{n_r}}\xrightarrow{\Pp}1.
\end{equation}
Consequently
\[
 \limsup_{n\in\mathcal N}\E\frac{Y_{(1)}(m_n,n)}n=\lambda.
\]
Along the selected subsequence $Y_{(1)}/n\to\lambda$ and
$(m_n-Y_{(1)})/n\to0$ both in probability and in $L^1$, and
$Y_{(2)}/n\to0$ in probability. For a full path, every fixed accuracy admits
whole intervals of consecutive sizes with uniformly vanishing failure
probability, as specified in Section~\ref{sec:paths}.
\end{corollary}

The zero-radius question is Problem~19.41 of \cite{Janson}, or
Problem~18.40 of \cite{JansonPreprint}. It asks whether $Y_{(1)}/n\to0$ in
probability can hold along the full sequence when $\lambda>\nu=0$.
For nontrivial weights with $w_0>0$, this is the zero-radius regime. Corollary~\ref{cor:paths} rules it
out without any restriction on the jumps or oscillations of $m_n$. This
is a subsequence conclusion, not a claim of condensation at every size or
of $\Pp(Y_{(1)}=m_n)\to1$. The source already supplies sparse-support
examples with sublinear maxima on selected subsequences.

The density condition on $\mathcal N$ in Corollary~\ref{cor:paths} is optimal
for a statement valid for every weight sequence: Section~\ref{sec:sharp}
gives an explicit counterexample whenever that condition fails.

Normalize $w_0=1$ by dividing all weights by $w_0$, which leaves every
allocation law unchanged. Let $d=\gcd\{j\ge1:w_j>0\}$. Dividing supported
sizes and admissible totals by $d$ preserves $Y_{(1)}/m$ and zero radius,
and replaces a density $\lambda$ by $\lambda/d$. Until the proof of
Theorem~\ref{thm:windows} is completed, we work with span one.

\section{Insertion without losing the small-box factor}\label{sec:insertion}
For an integer $u\ge0$ write
\[
 Z_u(t,n)=\sum_{\substack{x_1+\cdots+x_n=t\\0\le x_v\le u}}
             \prod_{v=1}^n w_{x_v}.
\]
The following exact finite inequality keeps the entire background partition
factor. It is applied only to positive-weight background count vectors.

\begin{lemma}[Occupation-count insertion]\label{lem:insertion}
Let $u,k,t,r,n$ be nonnegative integers, with $n\ge1$, $u<k$, $w_k>0$,
and $m=t+k+r$. Suppose $r$ has a fixed representation
by $R$ supported positive sizes at least $K$, together with $C$ other
supported positive sizes. All these sizes are less than $k$. Let $W_r$ be
the product of their weights. Put $Q=1+R+C$, and let $r_j$ be the
multiplicity of size $j$ in the representation. If $t+Q\le n$, then
\begin{equation}\label{eq:sharpinsertion}
 Z(m,n)\ge w_kW_r Z_u(t,n)
 \frac{\fall{n-t}{Q}}
 {\displaystyle\prod_{j:r_j>0}\fall{\min\{n,\lfloor m/j\rfloor\}}{r_j}}.
\end{equation}
In particular, if
\begin{equation}\label{eq:insertionconditions}
 nK\ge2m,\qquad t+1+R+C\le n/2,
\end{equation}
then
\begin{equation}\label{eq:insertion}
 Z(m,n)\ge\frac{n}{2^{C+1}}\,w_k W_r Z_u(t,n).
\end{equation}
Repeated sizes, including sizes occurring in both parts of the
representation, are allowed. An empty product equals one, and
$(x)_{\underline q}=x(x-1)\cdots(x-q+1)$ for $q\ge1$, with
$(x)_{\underline0}=1$.
\end{lemma}

\begin{proof}
Describe a background configuration by its occupation counts $a_j$, where
$\sum a_j=n$ and $\sum j a_j=t$, and $a_j=0$ for $j>u$. The total weight of
that count vector, including all coordinate permutations, is
\[
 \mathcal W(a)=\frac{n!}{\prod_j a_j!}\prod_j w_j^{a_j}.
\]
Set $Q=1+R+C$. Use $Q$ empty coordinates to insert one size-$k$ box and the
chosen representation of $r$. If $r_j$ is the multiplicity of size $j$ in
that representation, the resulting count vector is
\[
 b_0=a_0-Q,\quad b_k=1,\quad b_j=a_j+r_j\quad(0<j<k).
\]
The map $a\mapsto b$ is injective. Since $a_0\ge n-t$, its weight ratio is
\begin{equation}\label{eq:typeratio}
 \frac{\mathcal W(b)}{\mathcal W(a)}
 =w_kW_r\frac{\fall{a_0}{Q}}
 {\displaystyle\prod_{0<j<k}\prod_{\ell=1}^{r_j}(a_j+\ell)}.
\end{equation}
Since the final count $b_j=a_j+r_j$ is at most
$\min\{n,\lfloor m/j\rfloor\}$, bounding the numerator below by
$\fall{n-t}{Q}$ proves \eqref{eq:sharpinsertion}. Its denominator is
positive because $r_j\le\min\{n,\lfloor m/j\rfloor\}$, by the representation
and $Q\le n$. For the simpler bound, the numerator is at
least $(n/2)^Q$ by \eqref{eq:insertionconditions}.
Each denominator factor assigned to one of the $R$ large pieces is at most
$m/K$, since the final configuration has total $m$. Each of the other $C$
factors is at most $n$. When the two classes share a size, assign their
factors in either fixed order; both bounds still apply. Thus the ratio in
\eqref{eq:typeratio} is at least
\[
 w_kW_r\frac{n}{2^{C+1}}
       \left(\frac{nK}{2m}\right)^R
 \ge w_kW_r\frac{n}{2^{C+1}}.
\]
Summing over the distinct background count vectors gives
\eqref{eq:insertion}. Weights of sizes not used in the construction only
increase the original denominator. No assumption $m<2k$ is required:
every image count vector has exactly one size-$k$ box by construction.
\end{proof}

The sharp inequality can also be obtained directly with labelled boxes.
For every background, choose $Q$ ordered distinct empty coordinates and
insert the giant and a fixed ordered list of reservoir pieces. There are
at least $\fall{n-t}{Q}$ choices. A target has one possible giant location
and at most
$\prod_j\fall{\min\{n,\lfloor m/j\rfloor\}}{r_j}$ preimages of this kind.
The weight multiplier is always $w_kW_r$. Weighted double counting gives
\eqref{eq:sharpinsertion} without passing through occupation classes.

\section{Geometric records and two summation bounds}\label{sec:records}
Fix a real base $a>1$ and a constant $b>2$. For every sufficiently large
integer $i$ define
\begin{equation}\label{eq:records}
 B_i=\max_{\substack{1\le s\le\lfloor a^i\rfloor\\w_s>0}}
       \frac{\log w_s+b\log s}{s},\qquad
 s_i=\min\operatorname*{arg\,max}_{s}\frac{\log w_s+b\log s}{s},
 \qquad p_i=\frac{b\log s_i}{s_i}.
\end{equation}
Then $B_i$ is nondecreasing and tends to infinity. Also $s_i\to\infty$.
Set $d_i=B_i-B_{i-1}$ and call an index \emph{active} if $d_i>0$. At an
active index $J$,
\begin{equation}\label{eq:active}
 a^{J-1}<s_J\le a^J,\qquad
 \log w_{s_J}=s_JB_J-b\log s_J.
\end{equation}
At inactive indices the smallest maximizing index remains unchanged.
Consequently $s_i$ and $p_i$ are constant on each plateau of $B_i$.

\begin{lemma}[Summable plateau charges]\label{lem:plateau}
Fix a sufficiently large starting index $i_0$ and an integer $D\ge1$.
If $\mathcal A$ is the set of active indices after $i_0$, then
\begin{equation}\label{eq:plateausum}
 \sum_{J\in\mathcal A}\ \sum_{i=i_0}^{J-D} a^{i-J}p_i<\infty.
\end{equation}
More precisely, this sum is at most
\[
 \frac{a}{(a-1)^2}
 \left(p_{i_0}+\sum_{\ell\in\mathcal A}p_\ell\right),
\]
and the parenthesized sum is finite.
\end{lemma}

\begin{proof}
For an active $\ell$, \eqref{eq:active} gives
\[
 p_\ell\le ba\ell(\log a)a^{-\ell},
\]
which is summable. Consider a plateau starting at $\ell$ and ending just
before the next active index $v$. A term from this plateau can contribute
to \eqref{eq:plateausum} only at an active $J\ge v$: the starting active
index $J=\ell$ has cutoff $J-D<\ell$. Its total contribution is at most
\[
 p_\ell
 \left(\sum_{i=\ell}^{v-1}a^i\right)
 \left(\sum_{J=v}^{\infty}a^{-J}\right)
 \le \frac{a}{(a-1)^2}p_\ell.
\]
The same bound applies to the initial plateau at $i_0$. There are always
later active indices because $B_i\to\infty$. Summing proves the lemma.
\end{proof}

\begin{lemma}[Contractive memory kernel]\label{lem:kernel}
For $D\ge2$ put
\begin{equation}\label{eq:theta}
 C_a=2a^2,\qquad
 \Theta(a,D)=2a^{3-D}
 \left(\frac{D-1}{a-1}+\frac{1}{(a-1)^2}\right).
\end{equation}
For every nondecreasing sequence $B_i$ with the increments defined above,
and every finite upper limit $N$,
\begin{equation}\label{eq:kernelsum}
 C_a\sum_{\substack{i_0<J\le N\\J\ {\rm active}}}
       \sum_{i=i_0}^{J-D}a^{i-J}(B_{J-1}-B_i)
 \le\Theta(a,D)\sum_{\ell=i_0+1}^{N}d_\ell.
\end{equation}
In particular $D$ can be chosen so that $\Theta(a,D)<1/4$.
\end{lemma}

\begin{proof}
Expand $B_{J-1}-B_i=\sum_{\ell=i+1}^{J-1}d_\ell$. All terms are
nonnegative, so extend the sum over $J$ to all larger integers and the
lower limit on $i$ to $-\infty$ to bound each coefficient. For a fixed
$\ell$ and lag $t=J-\ell\ge1$, that coefficient, before multiplication by
$C_a$, is
\[
 \sum_{i\le\min(J-D,\ell-1)}a^{i-J}
 =\begin{cases}
  a^{1-D}/(a-1),&1\le t\le D-1,\\
  a^{-t}/(a-1),&t\ge D.
 \end{cases}
\]
Summing over $t$ gives
$a^{1-D}((D-1)/(a-1)+1/(a-1)^2)$. Multiplying by $2a^2$ gives
\eqref{eq:theta}. Exponential decay in $D$ proves the last assertion.
\end{proof}

\section{A supported representation of every small residual}\label{sec:residual}
The span-one reduction provides a finite set of positive supported sizes
whose gcd is one. Its nonnegative additive semigroup contains every integer
above some $h_*\ge1$. Here is a direct justification. Choose one member
$s$ of the finite set. The nonnegative residues generated modulo $s$ form
a subgroup of the finite cyclic group and hence all residues. Choose one
representable integer in each residue class; adding copies of $s$ then
represents every sufficiently large integer.

Fix an integer $K\ge2$. Enlarge $i_0$ so that, for all $i\ge i_0$,
\begin{equation}\label{eq:start}
 s_i\ge K,\qquad B_i\ge0,\qquad a^{i_0}\ge h_*.
\end{equation}
Choose once and for all a supported representation for each integer in
\[
 h_*\le v<2a^{i_0}.
\]
Let $C_*$ be a common bound on their numbers of positive parts, let
$M=\lceil2a^{i_0}\rceil$, and choose $C_0\ge0$ such that the logarithm of
each of their product weights is at least $-C_0$. These are finite constants.

\begin{lemma}[Greedy residual bound]\label{lem:greedy}
Suppose $J\ge i_0+D$, $r\ge h_*$, and
\[
 I(r):=\left\lfloor\log_a(r/2)\right\rfloor\le J-D,
 \qquad D\ge2.
\]
There is a representation of $r$ using supported parts at least $K$, and
at most $C_*$ additional supported parts, with product weight $W_r$ satisfying
\begin{equation}\label{eq:greedy}
 B_{J-1}r-\log W_r
 \le 2\sum_{i=i_0}^{J-D}a^{i+1}
       (B_{J-1}-B_i+p_i)+MB_{J-1}+C_0.
\end{equation}
All parts are at most $r$. The construction uses at most one stage at each
geometric index $i$.
\end{lemma}

\begin{proof}
With current residual $x$, stop if $x<2a^{i_0}$ and use its fixed terminal
representation. Otherwise set
\[
 i=\lfloor\log_a(x/2)\rfloor,\quad
 q=\left\lfloor\frac{x-h_*}{s_i}\right\rfloor,\quad
 x'=x-qs_i.
\]
Here $i\ge i_0$, $q\ge1$, and
\[
 h_*\le x'<h_*+s_i\le2a^i,\qquad qs_i\le x<2a^{i+1}.
\]
Thus the next geometric index is strictly smaller than $i$. Add $q$ copies
of $s_i$ and repeat. The procedure terminates in the fixed terminal range.
At stage $i$ the logarithmic weight per unit mass is $B_i-p_i$.
Since $B_{J-1}\ge B_i$ for every index used, that stage's deficit against
$B_{J-1}$ is at most
$2a^{i+1}(B_{J-1}-B_i+p_i)$. The terminal deficit is at most
$MB_{J-1}+C_0$. Unused summands are nonnegative, giving \eqref{eq:greedy}.
\end{proof}

\section{Deterministic margins}\label{sec:margin}
Fix $a>1$, $b>2$, and $D\ge2$ with $\Theta(a,D)<1/4$. Choose the finite
terminal data $i_0,M,C_0,C_*$ from the preceding section. For every active
$J\ge i_0+D$ set $k=s_J$, $h=B_{J-1}$, and define
\begin{align}
 E_J={}&2\sum_{i=i_0}^{J-D}a^{i+1}(B_{J-1}-B_i+p_i)
             +MB_{J-1}+C_0,\label{eq:energy}\\
 G_J={}&k(B_J-B_{J-1})-b\log k-E_J.\label{eq:margin}
\end{align}
All these quantities use finitely many weights. In particular $G_J$ does
not depend on a mass path or on a conditional allocation probability.

\begin{lemma}[Unbounded favorable margins]\label{lem:margin}
For the fixed data above,
\begin{equation}\label{eq:marginlimsup}
 \limsup_{\substack{J\to\infty\\J\ {\rm active}}}G_J=+\infty.
\end{equation}
\end{lemma}

\begin{proof}
Suppose instead $G_J\le T$ for all sufficiently large active $J$.
From \eqref{eq:margin}, $k>a^{J-1}$ and $\log k\le J\log a$, we obtain
\begin{align}
 d_J\le{}&C_a\sum_{i=i_0}^{J-D}a^{i-J}(B_{J-1}-B_i+p_i)\notag\\
 &+Ma^{1-J}B_{J-1}
    +a^{1-J}(bJ\log a+C_2),\label{eq:recurrence}
\end{align}
where $C_2=\max\{0,C_0+T\}$. Fix $J_0$ beyond this threshold with
\begin{equation}\label{eq:kappa}
 \kappa=\frac{Ma^{1-J_0}}{a-1}<\frac14.
\end{equation}
For $N\ge J_0$, put $S_N=B_N-B_{J_0-1}$ and sum
\eqref{eq:recurrence} over active $J\in[J_0,N]$. The memory term is at
most $\Theta(a,D)S_N$ plus a fixed finite-prefix contribution, by
Lemma~\ref{lem:kernel}. All $p_i$ terms have a finite total, by
Lemma~\ref{lem:plateau}; the last term is also summable.
The terminal term satisfies
\begin{align*}
 M\sum_{J=J_0}^N a^{1-J}B_{J-1}
 &\le\frac{MB_{J_0-1}a^{2-J_0}}{a-1}
    +M\sum_{\ell=J_0}^{N-1}d_\ell
                    \sum_{J=\ell+1}^{\infty}a^{1-J}\\
 &\le C_3+\kappa S_N.
\end{align*}
Thus $S_N\le(\Theta(a,D)+\kappa)S_N+C_4$, with $C_4$ independent of $N$.
The coefficient is below $1/2$, contradicting $B_N\to\infty$.
\end{proof}

The preceding proof yields a finite-prefix version, without assuming any
bound on $G_J$. It will also make the cancellation of all errors explicit.
For $L\ge0$ define
\[
 T_L=\sum_{j=L}^{\infty}a^{1-j}=\frac{a^{2-L}}{a-1},\qquad
 U_L=\sum_{j=L}^{\infty}j a^{1-j}
 =a^{2-L}\left(\frac{L}{a-1}+\frac1{(a-1)^2}\right).
\]
Set
\[
 \Pi=\frac{a}{(a-1)^2}
   \left(p_{i_0}+ba\log a\sum_{\ell=i_0+1}^{\infty}\ell a^{-\ell}\right)
\]
and, for $J_0\ge i_0+D$ with \eqref{eq:kappa}, set
\begin{align}
 \mathcal C={}&\Theta(a,D)(B_{J_0-1}-B_{i_0})+C_a\Pi
       +MB_{J_0-1}T_{J_0}\notag\\
       &+b(\log a)U_{J_0}+C_0T_{J_0}.\label{eq:explicitC}
\end{align}
These are finite nonnegative constants. The two infinite sums have displayed
geometric closed forms (the sum in $\Pi$ equals $U_{i_0+1}/a$).

\begin{lemma}[Finite-prefix margin bound]\label{lem:prefix}
For all $N\ge J_0$,
\begin{equation}\label{eq:finiteprefix}
 \sum_{\substack{J_0\le J\le N\\J\ {\rm active}}}\frac{G_J}{s_J}
 \ge(1-\Theta(a,D)-\kappa)(B_N-B_{J_0-1})-\mathcal C.
\end{equation}
In particular, with $x_+=\max\{x,0\}$ and the maximum of an empty set
interpreted as zero,
\begin{equation}\label{eq:prefixmax}
 \max_{\substack{J_0\le J\le N\\J\ {\rm active}}}(G_J)_+
 \ge\frac{(1-\Theta(a,D)-\kappa)(B_N-B_{J_0-1})-\mathcal C}{T_{J_0}}.
\end{equation}
\end{lemma}

\begin{proof}
Use the identity
$d_J=G_J/s_J+b\log(s_J)/s_J+E_J/s_J$ and sum over active indices.
The memory term and the terminal term have the bounds from the preceding
proof. Lemma~\ref{lem:plateau} and
$p_\ell\le ba\ell(\log a)a^{-\ell}$ at active $\ell$ bound the penalty
term by $C_a\Pi$. Finally,
\[
 \sum_{J\ge J_0}\frac{b\log s_J+C_0}{s_J}
 \le b(\log a)U_{J_0}+C_0T_{J_0},
\]
where sums on the left are restricted to active indices. Rearrangement
gives \eqref{eq:finiteprefix}. Since
$\sum_{J\ge J_0,\,J\ {\rm active}}s_J^{-1}\le T_{J_0}$, bounding each
$G_J$ by the maximum of their positive parts gives \eqref{eq:prefixmax}.
\end{proof}

\section{Uniform windows}\label{sec:windows}
Fix $\varepsilon\in(0,1)$ and $0<\alpha<\beta<\infty$. Choose $a>1$ close
to one, fix $b=4$, and choose $D\ge2$ with $\Theta(a,D)<1/4$. Next choose
$\delta>0$ such that
\begin{equation}\label{eq:parameters}
 \delta<\min\left\{\frac{a^{-D}}{100},\frac1{100(1+\beta)},\frac1{100}\right\},
 \qquad \frac1{a(1+5\delta)}>1-\varepsilon.
\end{equation}
Choose an integer $K\ge4(1+\beta)$, and then the fixed terminal data from
Section~\ref{sec:residual}. Every parameter is now fixed, before letting $J\to\infty$.
At an active $J$ write $k=s_J$, $u=\lfloor a^{J-1}\rfloor$, $h=B_{J-1}$,
and define the finite auxiliary distribution
\begin{equation}\label{eq:auxiliary}
 q_j=w_je^{-hj},\quad A_J=\sum_{j=0}^{u}q_j,\quad
 p_J(j)=q_j/A_J,\quad \mu_J=\sum_{j=0}^{u}j p_J(j).
\end{equation}
The record definition implies $q_0=1$ and $q_j\le j^{-b}$ for $1\le j\le u$.
Since $h\to\infty$, dominated convergence gives $A_J\to1$ and $\mu_J\to0$,
using $\sum j^{1-b}<\infty$ for the latter assertion.

\begin{lemma}[Uniform finite-size certificate]\label{lem:certificate}
For all sufficiently large active $J$, every integer pair with
\begin{equation}\label{eq:slab}
 (1+3\delta)k\le m\le(1+5\delta)k,\qquad
 \alpha\le m/n\le\beta
\end{equation}
is admissible in the span-one model, and
\begin{equation}\label{eq:certificate}
 \Pp_w(Y_{(1)}(m,n)\le u)\le\min\{1,C e^{-G_J}\},\qquad
 C=\max\left\{1,\frac{2^{C_*+3}\beta\delta}{1+3\delta}\right\}.
\end{equation}
Here $G_J$ is the deterministic quantity in \eqref{eq:margin}; the same
bound holds for every pair in the window.
\end{lemma}

\begin{proof}
The sufficiently large condition can be checked using only finite data:
require $J\ge i_0+D$, $\delta k\ge1$, $2\delta k\ge h_*$,
\begin{equation}\label{eq:finiteconditions}
 \mu_J\le\frac{\alpha\delta}{2(1+5\delta)},\qquad
 5\delta k+1\le\frac{(1+3\delta)k}{2\beta},
\end{equation}
and $h\ge0$. These all hold eventually. For the sum $S_n$ of $n$
independent copies of \eqref{eq:auxiliary}, Markov's inequality gives
$\Pp(S_n\le\delta k)\ge1/2$, uniformly in \eqref{eq:slab}, because
$n\le(1+5\delta)k/\alpha$. Hence some integer
$t\in[0,\lfloor\delta k\rfloor]$ satisfies
\begin{equation}\label{eq:atom}
 \Pp(S_n=t)\ge\frac1{2(\lfloor\delta k\rfloor+1)}.
\end{equation}
Put $r=m-k-t$, so $2\delta k\le r\le5\delta k$. By
\eqref{eq:parameters}, $I(r)\le J-D$. Lemma~\ref{lem:greedy} supplies a
representation with product $W_r$ and deficit $hr-\log W_r\le E_J$.
If it uses $R$ large pieces and $C'\le C_*$ terminal pieces, then
\[
 t+1+R+C'\le t+1+r=m-k+1\le5\delta k+1\le n/2.
\]
Also $nK\ge2m$, and every inserted part is smaller than $k$, since
$r\le5\delta k<k$. Lemma~\ref{lem:insertion} gives
\[
 Z(m,n)\ge\frac{n}{2^{C_*+1}}w_kW_r Z_u(t,n)>0.
\]
Thus no admissibility assumption has been used in producing the bound.
The exact identities and inequality
\[
 Z_u(t,n)=e^{ht}A_J^n\Pp(S_n=t),\qquad
 Z_u(m,n)=e^{hm}A_J^n\Pp(S_n=m)\le e^{hm}A_J^n
\]
now yield
\begin{align*}
 \Pp_w(Y_{(1)}\le u)
 &\le\frac{2^{C_*+2}(\lfloor\delta k\rfloor+1)}n
       \exp\{hk-\log w_k+hr-\log W_r\}\\
 &\le\frac{2^{C_*+3}\beta\delta}{1+3\delta}\,e^{-G_J}.
\end{align*}
The whole small-box factor $A_J^n$ cancels exactly. The last line uses
$\delta k\ge1$, $n\ge(1+3\delta)k/\beta$, and
$\log w_k=kB_J-b\log k$. This proves \eqref{eq:certificate}.
\end{proof}

\begin{proof}[Proof of Theorem~\ref{thm:windows}]
Lemma~\ref{lem:margin} supplies active $J_r\uparrow\infty$ with
$G_{J_r}\to\infty$. Let $k_r=s_{J_r}$. For all pairs in the corresponding
window,
\[
 \frac{u}{m}\ge\frac{a^{J_r-1}-1}{(1+5\delta)a^{J_r}}
 \longrightarrow\frac1{a(1+5\delta)}>1-\varepsilon
 \quad\text{as a lower bound}.
\]
Thus \eqref{eq:certificate} proves \eqref{eq:uniform}. The integer windows
have positive linear length, and $\alpha<\beta$ guarantees that for every
sufficiently large integer $m$ in such a window there is an integer $n$
with $\alpha\le m/n\le\beta$. Admissibility follows from the certificate.
For span $d$, apply the argument to $w_{dj}$, $m/d$ and the density interval
$[\alpha/d,\beta/d]$, then replace the supported sizes by $d k_r$.
The original windows contain multiples of $d$ for all large $r$, so they
also contain admissible pairs. The ratio $Y_{(1)}/m$ is unchanged.
\end{proof}

\section{Consequences for full paths}\label{sec:paths}
\begin{proof}[Proof of Corollary~\ref{cor:paths}]
Fix $\varepsilon>0$ and apply Theorem~\ref{thm:windows} with the density
interval $[\lambda/2,2\lambda]$. For each resulting $k_r$, let $n_r$ be the
first element of $\mathcal N$ at least $(1+4\delta)k_r/\lambda$.
The condition $N_{j+1}/N_j\to1$ implies
\[
 \frac{n_r}{k_r}\to\frac{1+4\delta}{\lambda},\qquad
 \frac{m_{n_r}}{k_r}\to1+4\delta.
\]
Thus these pairs eventually lie in the windows, and their failure
probabilities at accuracy $\varepsilon$ tend to zero. Apply this assertion
separately with $\varepsilon=1/r$, selecting a size greater than the preceding
one and a failure probability below $1/r$. This diagonal choice proves
\eqref{eq:conclusion}. The bounds $0\le Y_{(1)}/m_n\le1$ and
$m_n/n\to\lambda$ give the expectation and $L^1$ assertions. Finally
$Y_{(2)}\le m_n-Y_{(1)}$.
\end{proof}

For a full path $\mathcal N=\N$, the conclusion for a fixed accuracy is
stronger than the selection of one size per record. Define
\begin{equation}\label{eq:sizeintervals}
 I_r=\left[
 \left\lceil\frac{(1+\tfrac72\delta)k_r}{\lambda}\right\rceil,
 \left\lfloor\frac{(1+\tfrac92\delta)k_r}{\lambda}\right\rfloor
 \right]\cap\N.
\end{equation}
Uniformly for $n\in I_r$, the relation $m_n-\lambda n=o(n)$ implies
$|m_n-\lambda n|=o(k_r)$. Hence all such pairs eventually lie in
\eqref{eq:slab}, and
\begin{equation}\label{eq:blocks}
 \sup_{n\in I_r}\Pp_w(Y_{(1)}(m_n,n)\le(1-\varepsilon)m_n)\to0,
 \qquad |I_r|\sim\frac{\delta k_r}{\lambda}.
\end{equation}
In particular the intervals have a fixed positive relative width at every
fixed accuracy. They need not occur at regularly spaced locations.
No rate uniform in the weight sequence is asserted.

For simply generated plane trees the allocation has $m=n-1$. The cyclic
lemma groups every allocation sequence of total $n-1$ into an orbit of
$n$ rotations containing exactly one valid tree sequence; product weights
and the maximum are unchanged. Thus the maximum outdegree has the same
law as $Y_{(1)}$. Admissible sufficiently large tree sizes form the lattice
$n\equiv1\pmod d$: to verify eventual admissibility, choose a finite subset
of supported degrees of gcd $d$, represent every sufficiently large
multiple of $d$ as their nonnegative sum, and pad with zeros to $n$ entries.
There are at most $n-1$ positive summands since every degree is at least one.
The dense-index extension therefore applies.

\begin{corollary}[Subsequence concentration at the root]\label{cor:root}
For simply generated plane trees with zero-radius weights and $w_0>0$,
there are admissible $n_r\uparrow\infty$ such that the root has
$(1-o_{\Pp}(1))n_r$ children. With probability tending to one the root is
the unique vertex of maximum outdegree. The number of vertices at depth
at least two is $o_{\Pp}(n_r)$, also in expectation after division by $n_r$.
\end{corollary}

\begin{proof}
Fix $n\ge2$. Let $A_j$ be the number of vertices of outdegree $j$, and let
$D=\max\{j:A_j>0\}$. Conditional on $(A_j)$, the known root-degree identity
\cite[Lemma~14.7]{JansonPreprint} reads
\begin{equation}\label{eq:rootidentity}
 \Pp(d^+(o)=j\mid(A_i))=\frac{jA_j}{n-1}.
\end{equation}
For a direct count, the number of plane trees with these degree counts is
$(n-1)!/\prod_iA_i!$. Removing a root of degree $j$ leaves an ordered forest
of $j$ trees. The cyclic lemma counts these forests as
$j(n-2)!/((A_j-1)!\prod_{i\ne j}A_i!)$, yielding
\eqref{eq:rootidentity}. All trees with given counts have the same product
weight, so conditioning on the counts does not alter this calculation.

Along the subsequence from Corollary~\ref{cor:paths}, $D/(n-1)\to1$ in
probability. On $D>(n-1)/2$ the maximum is unique, and its probability of
being the root, conditionally on the counts, is $D/(n-1)$. Taking
expectations gives probability tending to one: explicitly, the error is
at most
$\Pp(D\le(n-1)/2)+\E[1-D/(n-1)]$, which tends to zero. This also proves
$d^+(o)/(n-1)\to1$ in probability and in $L^1$. Finally the number of
vertices at depth at least two is exactly $n-1-d^+(o)$.
\end{proof}

The root identity is not a new result. The corollary is its application
to the unrestricted subsequence statement proved above; it is stronger
than merely identifying an infinite-star local limit.

\section{The sampling condition is sharp}\label{sec:sharp}
\begin{theorem}[Sharp universal sampling criterion]\label{thm:sampling}
For an infinite index set $\mathcal N=\{N_1<N_2<\cdots\}$, the following
are equivalent:
\begin{enumerate}[label=(\roman*)]
\item $N_{j+1}/N_j\to1$.
\item For every fixed zero-radius weight sequence with $w_0>0$ and every
      eventually admissible mass path on $\mathcal N$ with positive finite
      limiting density, some subsequence satisfies $Y_{(1)}/m\to1$ in
      probability.
\end{enumerate}
If $\limsup_jN_{j+1}/N_j=\infty$, there are such weights, with $w_1>0$, for
which $Y_{(1)}(n-1,n)/n\to0$ deterministically throughout $n\in\mathcal N$.
\end{theorem}

\begin{proof}
The implication (i)$\Rightarrow$(ii) is Corollary~\ref{cor:paths}.
If (i) fails, there exist $\theta>1$ and infinitely many increasing indices
$r_j$ with $N_{r_j+1}\ge\theta N_{r_j}$. Discard finitely many so that
$N_{r_j}\ge2$. Set
\[
 S=\{0,1\}\cup\{N_{r_j}:j\ge1\},\qquad
 w_0=w_1=1,\quad w_k=e^{k^2}\ (k\in S\setminus\{0,1\}),
\]
and set every other weight to zero. These are fixed weights with zero
radius. The tree path $m_n=n-1$ is admissible at every size by using ones
and zeros. If a supported $k=N_{r_j}$ is available at a sampled size
$n\in\mathcal N$, then $k\le n-1$, so
$n\ge N_{r_j+1}\ge\theta k$. Consequently
\[
 Y_{(1)}(n-1,n)\le\max\{1,n/\theta\}
 \quad(n\in\mathcal N)
\]
for every configuration. This implies
$\limsup_{n\in\mathcal N}Y_{(1)}/(n-1)\le1/\theta<1$, disproving (ii).

For the final assertion, select $r_j$ so that
$N_{r_j+1}/N_{r_j}\ge j+1$ and use the same construction. Fix $L$.
Any available supported $N_{r_j}$ with $j\ge L$ is at most $n/(L+1)$.
The finitely many earlier supported values, divided by $n$, tend to zero.
Thus $\limsup_{n\in\mathcal N}Y_{(1)}/n\le1/(L+1)$ deterministically.
Let $L\to\infty$.
\end{proof}

This is an optimality statement about a universal theorem, not an assertion
that every weight sequence fails on a nondense set. For example, the powers
of two as sampled sizes and as nonzero large support already force a
maximum at most $n/2$ on the tree path; factorial gaps can force a vanishing
fraction, in the same spirit as Example~19.39 of \cite{Janson}.

\section{Exact finite illustrations}\label{sec:finite}
The argument uses a deterministic margin to give a uniform window statement
and a sharp sampling criterion. The finite checks below test its algebraic
components; they are not a proof-assistant formalization or an independent
audit of the limiting argument.

The accompanying \texttt{code/check\_uniform\_windows.py} verifies finite
occupation-count bounds and conditional probabilities on an entire explicit
window, the exact finite-prefix inequality with rational geometric data,
and the support obstructions at powers-of-two and factorial sample sizes.
It also executes the accompanying checks of the insertion, kernel, and
residual estimates. In the explicit window the
weights are supported on $\{0,2,3,5,17\}$ below 120, with
$w_j=2^{j^2}$ at the positive supported values; assigning the same weights
at $j=\ell!$ for $\ell\ge5$ makes this one fixed infinite zero-radius law.
Every pair $20\le m\le22$, $1/2\le m/n\le2$ is checked with exact arithmetic.
There are 97 such pairs. The sharp insertion bound certifies
\[
 \sup_{\substack{20\le m\le22\\1/2\le m/n\le2}}
 \Pp_w(Y_{(1)}\le3m/4)\le2^{-189}.
\]
The comparison is summarized in Table~\ref{tab:finite}. All probabilities
and bounds are computed as rational numbers; decimal values are displays,
not the basis of the assertions.
\begin{table}[ht]
\centering
\caption{An exact finite illustration over all 97 admissible pairs
$20\le m\le22$, $1/2\le m/n\le2$.}\label{tab:finite}
\begin{tabular}{@{}ll@{}}
\toprule
Quantity & Maximum over the window \\
\midrule
Actual failure probability & $1.50986220045\cdot10^{-58}$ \\
Sharp insertion upper bound & $9.05917320269\cdot10^{-58}\le2^{-189}$ \\
Coarse insertion upper bound & $7.58157669219\cdot10^{-56}\le2^{-183}$ \\
\bottomrule
\end{tabular}
\end{table}
The exact rational values are saved in the checker output. The
checker also verifies the root-degree identity by enumerating all plane
trees of sizes two through ten.
These finite computations do not replace the asymptotic argument.

The positive-boundary-mean question is outside the scope of the present
results. The factorial and successive-ratio results in \cite{JJS}, and the
size-biased root identity, are used only with their stated attribution.


\section*{Acknowledgements and funding}
The research presented in this paper was supported by the European Research
Council (ERC) under the European Union's Horizon Europe research and
innovation programme (grant agreement No.~101041711), by the Simons
Foundation as part of the Collaboration on the Mathematical and Scientific
Foundations of Deep Learning, by Heights Labs, by Convex Nexus Capital, by
the Israel Science Foundation (grant number 2258/19), and by the Israel
Science Foundation (ISF Grant 4101/25).

\section*{Declaration of competing interest}
The author holds equity in Heights Labs and in Convex Nexus Capital and
has a management role at Convex Nexus Capital. These entities are included
in the funding declaration above.

\section*{Declaration of generative AI and AI-assisted technologies}
This manuscript was developed with substantial assistance from OpenAI's
ChatGPT as part of an ongoing AI-assisted mathematical research project,
provisionally called Math Brain. AI assistance contributed to the
mathematical arguments, literature searches, computations, and writing.
Five documented review passes were performed by the same AI assistant;
they are not independent referee reports, human peer review, or a formal
verification of the asymptotic theorem. The author directed the project
and is responsible for decisions concerning dissemination. The present
version is a research draft for further mathematical review.

\section*{Data and code availability}
The complete \LaTeX{} source, exact-arithmetic Python scripts, saved outputs,
and five review reports accompany the manuscript. The scripts use only the
Python standard library. No external empirical dataset is required.
The development archive is maintained in the author's private repository
\texttt{yspennstate/math-open-problems}; public repository access is not
asserted. The auxiliary computations verify finite identities and
inequalities, not the asymptotic theorem by exhaustive testing.

\begin{thebibliography}{9}
\bibitem{Janson}
S. Janson,
\emph{Simply generated trees, conditioned Galton--Watson trees, random
allocations and condensation},
Probability Surveys \textbf{9} (2012), 103--252.
\href{https://doi.org/10.1214/11-PS188}{doi:10.1214/11-PS188}.
\bibitem{JansonPreprint}
S. Janson,
\emph{Simply generated trees, conditioned Galton--Watson trees, random
allocations and condensation}, preprint, December 2011.
\href{https://arxiv.org/abs/1112.0510v1}{arXiv:1112.0510v1}.
Numbering differs from the published version \cite{Janson}.
\bibitem{JJS}
S. Janson, T. Jonsson and S. \"O. Stef\'ansson,
\emph{Random trees with superexponential branching weights},
Journal of Physics A: Mathematical and Theoretical \textbf{44} (2011), 485002.
\href{https://arxiv.org/abs/1104.2810}{arXiv:1104.2810}.
\end{thebibliography}
\end{document}
